% GENERATED FILE. Edit canonical lessons, then run npm run generate:books.
% canonical-source-hash: fnv1a64:2825ce73b145ad25
% canonical-lessons: HI-C485-badhazmi, HI-C485-murcha, HI-C485-pecis, HI-C485-ubakna, HI-C485-dakar

\chapter{Feeling Unwell}
\label{ch:hi-feeling-unwell}
\hlchaptermodality{485}

\begin{chapteropening}
\textbf{By the end of this chapter:} \emph{I can say that someone felt dizzy or sick.}

Complete the last lesson of chapter 485: I can say that someone felt dizzy or sick.
\end{chapteropening}

\section[\texorpdfstring{\hi{बदहज़मी} (badhazmī)}{बदहज़मी (badhazmī)}]{\hi{बदहज़मी} (badhazmī) — indigestion}

\label{lesson:HI-C485-badhazmi}



\begin{quote}
Before the new one: say the Hindi for to rub in, then the Hindi for a razor.
\end{quote}

\subsection*{You'll want to know: \hi{बदहज़मी}}
\textbf{\hi{बदहज़मी}} — \emph{badhazmī} — \textquotedblleft{}indigestion\textquotedblright{}.

\begin{quote}
\hi{ज़्यादा} \hi{मिठाई} \hi{खाकर} \hi{अरुण} \hi{को} \hi{बदहज़मी} \hi{हो} \hi{गई।}
\end{quote}

\begin{quote}
\emph{zyādā miṭhāī khākar Aruṇ ko badhazmī ho gaī.}
\end{quote}

\begin{quote}
Arun got indigestion from eating too many sweets.
\end{quote}

Feminine.

\begin{cousinweb}
\textbf{\hi{बदहज़मी}} (\emph{badhazmī}) is usually traced to Persian \emph{bad-hazmī}, \textquotedblleft{}indigestion\textquotedblright{}, which joins Persian \emph{bad}, \textquotedblleft{}bad\textquotedblright{}, to Arabic \emph{haḍm}, \textquotedblleft{}digestion\textquotedblright{}, said \emph{hazm} in Persian.
\end{cousinweb}

\begin{grammarlens}[title={What you've built}]
One more word for describing what you lived and felt.
\end{grammarlens}

\subsection*{Your turn}
\begin{itemize}
\raggedright
  \item \emph{Say it:} \emph{badhazmī}
  \item \emph{Say it:} \emph{badhazmī}, once more
  \item \emph{Say it:} \emph{zyādā miṭhāī khākar Aruṇ ko badhazmī ho gaī.} — Arun got indigestion from eating too many sweets.
  \item \emph{Recall:} say the Hindi for to rub in, then the Hindi for a razor, then say \emph{badhazmī} again
\end{itemize}

\begin{tcolorbox}[breakable,colback=teal!4,colframe=teal!35!black,title={Before you move on}]
What does \hi{बदहज़मी} mean? (\textquotedblleft{}Indigestion\textquotedblright{}.) Say it once more.
\end{tcolorbox}
\section[\texorpdfstring{\hi{मूर्छा} (mūrchā)}{मूर्छा (mūrchā)}]{\hi{मूर्छा} (mūrchā) — a faint; a swoon}

\label{lesson:HI-C485-murcha}



\begin{quote}
Before the new one: say the Hindi for a razor, then the Hindi for indigestion.
\end{quote}

\subsection*{You'll want to know: \hi{मूर्छा}}
\textbf{\hi{मूर्छा}} — \emph{mūrchā} — \textquotedblleft{}a faint; a swoon\textquotedblright{}.

\begin{quote}
\hi{गर्मी} \hi{से} \hi{एक} \hi{आदमी} \hi{को} \hi{सड़क} \hi{पर} \hi{मूर्छा} \hi{आ} \hi{गई।}
\end{quote}

\begin{quote}
\emph{garmī se ek ādmī ko saṛak par mūrchā ā gaī.}
\end{quote}

\begin{quote}
A man fainted in the street from the heat.
\end{quote}

Feminine.

\begin{cousinweb}
\textbf{\hi{मूर्छा}} (\emph{mūrchā}) is a tatsama, borrowed from Sanskrit \emph{mūrchā}, \textquotedblleft{}fainting, a swoon, stupor\textquotedblright{} (Monier-Williams). \emph{mūrchā ānā} is to faint; everyday speech often says \emph{behoś honā}.
\end{cousinweb}

\begin{grammarlens}[title={What you've built}]
One more word for describing what you lived and felt.
\end{grammarlens}

\subsection*{Your turn}
\begin{itemize}
\raggedright
  \item \emph{Say it:} \emph{mūrchā}
  \item \emph{Say it:} \emph{mūrchā}, once more
  \item \emph{Say it:} \emph{garmī se ek ādmī ko saṛak par mūrchā ā gaī.} — A man fainted in the street from the heat.
  \item \emph{Recall:} say the Hindi for a razor, then the Hindi for indigestion, then say \emph{mūrchā} again
\end{itemize}

\begin{tcolorbox}[breakable,colback=teal!4,colframe=teal!35!black,title={Before you move on}]
What does \hi{मूर्छा} mean? (\textquotedblleft{}A faint; a swoon\textquotedblright{}.) Say it once more.
\end{tcolorbox}
\section[\texorpdfstring{\hi{पेचिश} (peciś)}{पेचिश (peciś)}]{\hi{पेचिश} (peciś) — dysentery}

\label{lesson:HI-C485-pecis}



\begin{quote}
Before the new one: say the Hindi for indigestion, then the Hindi for a faint.
\end{quote}

\subsection*{You'll want to know: \hi{पेचिश}}
\textbf{\hi{पेचिश}} — \emph{peciś} — \textquotedblleft{}dysentery\textquotedblright{}.

\begin{quote}
\hi{गंदा} \hi{पानी} \hi{पीकर} \hi{अरुण} \hi{को} \hi{पेचिश} \hi{हो} \hi{गई।}
\end{quote}

\begin{quote}
\emph{gandā pānī pīkar Aruṇ ko peciś ho gaī.}
\end{quote}

\begin{quote}
Arun got dysentery from drinking dirty water.
\end{quote}

Feminine.

\begin{cousinweb}
\textbf{\hi{पेचिश}} (\emph{peciś}) is usually traced to Persian \emph{pēchish}, \textquotedblleft{}twisting, gripes\textquotedblright{}, from \emph{pēch}, \textquotedblleft{}a twist\textquotedblright{}: the illness is named for the cramps in the belly.
\end{cousinweb}

\begin{grammarlens}[title={What you've built}]
One more word for describing what you lived and felt.
\end{grammarlens}

\subsection*{Your turn}
\begin{itemize}
\raggedright
  \item \emph{Say it:} \emph{peciś}
  \item \emph{Say it:} \emph{peciś}, once more
  \item \emph{Say it:} \emph{gandā pānī pīkar Aruṇ ko peciś ho gaī.} — Arun got dysentery from drinking dirty water.
  \item \emph{Recall:} say the Hindi for indigestion, then the Hindi for a faint, then say \emph{peciś} again
\end{itemize}

\begin{tcolorbox}[breakable,colback=teal!4,colframe=teal!35!black,title={Before you move on}]
What does \hi{पेचिश} mean? (\textquotedblleft{}Dysentery\textquotedblright{}.) Say it once more.
\end{tcolorbox}
\section[\texorpdfstring{\hi{उबकना} (ubaknā)}{उबकना (ubaknā)}]{\hi{उबकना} (ubaknā) — to retch; to throw up}

\label{lesson:HI-C485-ubakna}



\begin{quote}
Before the new one: say the Hindi for a faint, then the Hindi for dysentery.
\end{quote}

\subsection*{You'll want to know: \hi{उबकना}}
\textbf{\hi{उबकना}} — \emph{ubaknā} — \textquotedblleft{}to retch; to throw up\textquotedblright{}.

\begin{quote}
\hi{बासी} \hi{दूध} \hi{पीकर} \hi{बच्चा} \hi{उबक} \hi{गया।}
\end{quote}

\begin{quote}
\emph{bāsī dūdh pīkar baccā ubak gayā.}
\end{quote}

\begin{quote}
The child threw up after drinking stale milk.
\end{quote}

Past (he …): \hi{वह} \hi{उबक} \hi{गया}, \emph{vah ubak gayā}.

\begin{cousinweb}
\textbf{\hi{उबकना}} (\emph{ubaknā}) is a tadbhava: CDIAL lists Hindi \emph{ubaknā}, \textquotedblleft{}to vomit\textquotedblright{}, under \emph{*ubbakka}, a form it reconstructs (CDIAL 2337). \emph{ulṭī}, \textquotedblleft{}vomiting\textquotedblright{}, is the noun most people use.
\end{cousinweb}

\begin{grammarlens}[title={What you've built}]
One more word for describing what you lived and felt.
\end{grammarlens}

\subsection*{Your turn}
\begin{itemize}
\raggedright
  \item \emph{Say it:} \emph{ubaknā}
  \item \emph{Say it:} \emph{ubaknā}, once more
  \item \emph{Say it:} \emph{bāsī dūdh pīkar baccā ubak gayā.} — The child threw up after drinking stale milk.
  \item \emph{Recall:} say the Hindi for a faint, then the Hindi for dysentery, then say \emph{ubaknā} again
\end{itemize}

\begin{tcolorbox}[breakable,colback=teal!4,colframe=teal!35!black,title={Before you move on}]
What does \hi{उबकना} mean? (\textquotedblleft{}To retch; to throw up\textquotedblright{}.) Say it once more.
\end{tcolorbox}
\section[\texorpdfstring{\hi{डकार} (ḍakār)}{डकार (ḍakār)}]{\hi{डकार} (ḍakār) — a burp; a belch}

\label{lesson:HI-C485-dakar}



\begin{quote}
Before the new one: say the Hindi for dysentery, then the Hindi for to retch.
\end{quote}

\subsection*{You'll want to know: \hi{डकार}}
\textbf{\hi{डकार}} — \emph{ḍakār} — \textquotedblleft{}a burp; a belch\textquotedblright{}.

\begin{quote}
\hi{खाना} \hi{खाकर} \hi{दादा} \hi{ने} \hi{ज़ोर} \hi{से} \hi{डकार} \hi{ली।}
\end{quote}

\begin{quote}
\emph{khānā khākar dādā ne zor se ḍakār lī.}
\end{quote}

\begin{quote}
After eating, my grandfather gave a loud burp.
\end{quote}

Feminine.

\begin{cousinweb}
\textbf{\hi{डकार}} (\emph{ḍakār}) is a tadbhava: CDIAL lists Hindi \emph{ḍakār}, \textquotedblleft{}belch\textquotedblright{}, under \emph{*ḍakkāra}, a form it reconstructs (CDIAL 5521). \emph{ḍakār lenā} is to burp.
\end{cousinweb}

\begin{grammarlens}[title={What you've built}]
One more word for describing what you lived and felt.
\end{grammarlens}

\subsection*{Your turn}
\begin{itemize}
\raggedright
  \item \emph{Say it:} \emph{ḍakār}
  \item \emph{Say it:} \emph{ḍakār}, once more
  \item \emph{Say it:} \emph{khānā khākar dādā ne zor se ḍakār lī.} — After eating, my grandfather gave a loud burp.
  \item \emph{Recall:} say the Hindi for dysentery, then the Hindi for to retch, then say \emph{ḍakār} again
\end{itemize}

\begin{tcolorbox}[breakable,colback=teal!4,colframe=teal!35!black,title={Before you move on}]
What does \hi{डकार} mean? (\textquotedblleft{}A burp; a belch\textquotedblright{}.) Say it once more.
\end{tcolorbox}
